\documentclass[11pt,a4paper]{article}
\usepackage[utf8]{inputenc}
\usepackage[T1]{fontenc}
\usepackage[english]{babel}
\usepackage{lmodern}
\usepackage{geometry}
\usepackage{xcolor}
\usepackage{tocloft}
\usepackage{hyperref}
\usepackage{enumitem}
\usepackage{parskip}
\usepackage{booktabs}
\usepackage{longtable}
\usepackage{fancyhdr}
\usepackage{lastpage}
\usepackage{listings}
\usepackage{titlesec}
\geometry{margin=2.15cm}
\setlength{\headheight}{14pt}
\setlength{\emergencystretch}{3em}
\setlist{itemsep=3pt,parsep=1pt,topsep=4pt}
\definecolor{titleblue}{HTML}{123B5D}
\definecolor{softblue}{HTML}{EDF5FB}
\definecolor{softgray}{HTML}{F5F7FA}
\definecolor{accent}{HTML}{0F6E8C}
\hypersetup{
  colorlinks=true,linkcolor=titleblue,urlcolor=accent,
  pdftitle={Scientific Data Analysis v2.7: English publication edition},
  pdfauthor={interemi},
  pdfsubject={Senior integration hardening and preparation for Scientific Workbench synchronization}
}
\titleformat{\section}{\Large\bfseries\color{titleblue}}{\thesection.}{0.6em}{}
\titleformat{\subsection}{\large\bfseries\color{accent}}{\thesubsection.}{0.6em}{}
\setlength{\cftsecnumwidth}{2.8em}
\tocloftpagestyle{fancy}
\lstdefinestyle{promptstyle}{
  basicstyle=\ttfamily\small,backgroundcolor=\color{softgray},
  frame=single,rulecolor=\color{softblue},breaklines=true,columns=fullflexible
}
\pagestyle{fancy}
\fancyhf{}
\fancyhead[L]{Scientific Data Analysis}
\fancyhead[R]{v2.7 guide: English edition}
\fancyfoot[C]{\thepage\ of \pageref{LastPage}}
\newcommand{\code}[1]{\texttt{\detokenize{#1}}}

\begin{document}
\pagenumbering{gobble}
\begin{titlepage}
\centering
\vspace*{1.0cm}
{\Huge\bfseries\color{titleblue} Senior Integration Hardening\\[0.15cm]
Scientific Data Analysis v2.7\par}
\vspace{0.65cm}
{\Large Mother and children prepared for deliberate synchronization with Scientific Workbench\par}
\vspace{0.8cm}
{\large Architectural health: canonical registry, multi-root ownership,
accurate coverage claims, stub integrity, and an app-side plan\par}
\vspace{0.8cm}
{\large English publication edition: September 22, 2026\par}
\vspace{0.5cm}
\begin{minipage}{0.93\textwidth}
\small
\raggedright
Translation of the preserved private v2.7 guide at commit \code{d6583f1}.
Original TeX SHA-256:
\nolinkurl{9e10ef1a7e63f5e3f3ac6e41a0eb1d6d419f363fd8faf0dbb827ec10aef89c93}.
Private source:
\nolinkurl{skills/scientific-data-analysis/.cache/active_references/guides/guia_scientific_data_analysis_v2_7.tex}.

This is dated evidence, not current release verification. The original guide
prepares an integration handoff; it does not perform synchronization. Its
recorded ratings and open app-side work retain that historical meaning.
The original remains unchanged privately.
\end{minipage}
\vfill
{\large Prepared with Codex\par}
{\large Copyright 2026 interemi\par}
\end{titlepage}

\clearpage
\pagenumbering{roman}
\fancyfoot[C]{\thepage}
\begin{abstract}
\code{scientific-data-analysis} v2.7 continues v2.6 with a different focus:
addressing correctable risks identified by the senior AI / Agents / Plugins /
Skills evaluation before resynchronizing the skill family with Scientific
Workbench, rather than measuring local cost.

Mother and children already have a recorded \code{100/A} \code{plugin-eval}
rating. v2.7 adds the missing application contract: resolve compact canonical
registries, treat mother and children as distinct roots, describe smoke/regression
coverage accurately, verify symlink/stub integrity, and provide a short app-side
synchronization plan.
\end{abstract}
\tableofcontents
\newpage
\pagenumbering{arabic}
\fancyfoot[C]{\thepage\ of \pageref{LastPage}}

\section{Executive overview}

v2.7 is a senior hardening release. It creates no new capabilities, changes no
scientific contracts, and does not edit Scientific Workbench. It prepares the
mother and child skills for app consumption without fragile heuristics.

The product decisions are:
\begin{itemize}[leftmargin=1.5em]
\item \code{scientific-data-analysis} remains the normal entrypoint for broad tasks.
\item \code{scientific-data-astro}, \code{scientific-data-documents}, \code{scientific-data-notebooks}, and \code{scientific-data-maintainer} are independent roots for app-side consumption.
\item Scientific Workbench must resolve \code{canonical_registry} recursively.
\item Every capability must execute from its owning root, rather than always from the mother.
\item Current coverage is smoke/regression coverage, not exhaustive coverage.
\item Maintainer remains maintainer-only and must not appear as an ordinary action.
\end{itemize}

\section{From v2.6 to v2.7}
\begin{longtable}{p{0.31\linewidth}p{0.57\linewidth}}
\toprule
\textbf{Area} & \textbf{v2.7 change} \\
\midrule
Senior evaluation & Uses the senior evaluation dated July 1, 2026. \\
P1 app-side & Documents canonical registry resolution and multi-root ownership. \\
P2 coverage & Documents that \code{coverage-summary.json} describes smoke/regression coverage, not exhaustive coverage. \\
P2 integrity & Adds regression checks for broken symlinks, unresolved stubs, and incomplete canonical registries. \\
P2 drift & Freezes the app-side handoff to reduce drift between exposure/app-readiness and Swift code. \\
Concise documentation & Adds \code{v2-7-scientificworkbench-sync-plan.md} as a short guide for the app task. \\
Guide & This guide replaces v2.6 as the active document. \\
\bottomrule
\end{longtable}

\section{Applied corrections}
\begin{longtable}{p{0.27\linewidth}p{0.18\linewidth}p{0.43\linewidth}}
\toprule
\textbf{Senior finding} & \textbf{Status} & \textbf{Correction} \\
\midrule
Resolve \code{canonical_registry} & PASS & Recursive regression reaches YAML with 61 entries. \\
Do not assume one root & PASS skill-side / WARNING app-side & Multi-root contract documented; app synchronization remains pending. \\
Non-exhaustive coverage & PASS & Artifacts declare smoke/regression coverage, not exhaustive coverage. \\
Symlinks/stubs/wrappers & PASS & Regression checks for broken symlinks in \code{references/}. \\
Exposure/app-readiness drift & Documentation PASS & Handoff requires app drift tests before synchronization is declared. \\
Dense documentation & PASS & Short plan dedicated to Scientific Workbench. \\
\bottomrule
\end{longtable}

\section{Multi-root contract for Scientific Workbench}

The app must load these five installed roots:
\begin{longtable}{p{0.24\linewidth}p{0.58\linewidth}}
\toprule
\textbf{Owner} & \textbf{Installed root} \\
\midrule
mother & \code{~/.codex/skills/scientific-data-analysis} \\
astro & \code{~/.codex/skills/scientific-data-astro} \\
documents & \code{~/.codex/skills/scientific-data-documents} \\
notebooks & \code{~/.codex/skills/scientific-data-notebooks} \\
maintainer & \code{~/.codex/skills/scientific-data-maintainer} \\
\bottomrule
\end{longtable}

The mother may remain the visible entrypoint for broad tasks, but command
construction must use each capability's owning root as \code{workingDirectory}.
This prevents the app from consuming an incomplete or outdated system view.

\section{Canonical registry contract}

Each \code{public_surface_registry.yaml} may be compact. A consumer must:
\begin{enumerate}[leftmargin=1.7em]
\item open the root's registry;
\item resolve \code{canonical_registry} relative to the stub when present;
\item repeat until reaching YAML with \code{entries:};
\item verify the expected 61-entry surface;
\item retain a fingerprint/cache without rewriting the skill.
\end{enumerate}

\section{v2.7 regression}

New regression:
\begin{lstlisting}[style=promptstyle]
python scripts/audit_v2_7_senior_integration_regression.py
\end{lstlisting}

It should return \code{PASS} with an accepted warning if Scientific Workbench
still uses a single \code{DefaultPaths.skillRoot}. That warning is not a skill
failure; it identifies the work to complete during app synchronization.

\section{Recommended gates}
\begin{lstlisting}[style=promptstyle]
python3 -m py_compile scripts/audit_v2_7_senior_integration_regression.py
python scripts/audit_v2_7_senior_integration_regression.py
python scripts/audit_v2_7_plugin_eval_coverage_artifacts.py
python scripts/audit_v2_6_benchmark_matrix_regression.py
python scripts/audit_v2_5_budget_deferral_regression.py
python scripts/sync_public_surface_docs.py --check
python scripts/portable_smoke_test.py --profile core
\end{lstlisting}

In maintainer, the complete regression lives in \code{fixtures/}; the mother
retains a compatible wrapper.

\newpage
\section{Short Scientific Workbench synchronization plan}

App-side synchronization must follow this order:
\begin{enumerate}[leftmargin=1.7em]
\item Verify the installed skill using the v2.7 gates.
\item Create \code{SkillRootCatalog} or its equivalent in the app.
\item Resolve \code{canonical_registry} recursively for each root.
\item Associate every \code{CapabilityEntry} with its owning root.
\item Update \code{CapabilityCommandBuilder} to execute from the owning root.
\item Keep \code{ToolEnvelopeParser} and \code{ArtifactDiscovery} compatible with v1.8, v1.9, and v2.x.
\item Hide maintainer-only actions in normal mode.
\item Run Swift tests, the quality gate, and smoke cases for tables, documents, notebooks, FITS, and an absent optional backend.
\end{enumerate}

\section{Non-goals}
\begin{itemize}[leftmargin=1.5em]
\item Do not synchronize Scientific Workbench during this phase.
\item Do not copy skills into the app bundle.
\item Do not turn maintainer into an ordinary user action.
\item Do not present smoke/regression coverage as exhaustive.
\item Do not change correct scientific behavior to satisfy the app.
\end{itemize}

\section{Conclusion}

v2.7 prepares the skill family with a recorded automatic \code{100/A} rating,
approved senior evaluation, closed skill-side P1/P2 corrections, and a concrete
plan for Scientific Workbench to stop depending on a single root. The recorded
next step is deliberate application synchronization using this contract.

\end{document}
